% Shared preamble for the literature/survey LaTeX documents.
% Used as the very first line of each document:  % Shared preamble for the literature/survey LaTeX documents.
% Used as the very first line of each document:  % Shared preamble for the literature/survey LaTeX documents.
% Used as the very first line of each document:  \input{preamble}
% (\input is a TeX primitive, so it is legal before \begin{document}.)
\documentclass[11pt]{article}
\usepackage[T1]{fontenc}
\usepackage{lmodern}
\usepackage[letterpaper,margin=1in]{geometry}
\usepackage{amsmath,amssymb}
\usepackage{array,booktabs}
\usepackage{enumitem}
\usepackage[most]{tcolorbox}
\usepackage{xcolor}
\usepackage[colorlinks=true,linkcolor=teal!60!black,urlcolor=teal!60!black,citecolor=teal!60!black]{hyperref}
\usepackage[backend=biber,style=numeric-comp,sorting=ynt,maxbibnames=8,maxcitenames=2,giveninits=true]{biblatex}
\addbibresource{references.bib}
\newtcolorbox{callout}{colback=black!3,colframe=black!25,boxrule=0.4pt,arc=2pt,
  left=6pt,right=6pt,top=4pt,bottom=4pt,breakable}
\setlist{topsep=3pt,itemsep=1pt,parsep=1pt,leftmargin=1.4em}
\setlength{\parindent}{0pt}\setlength{\parskip}{0.5em}
\sloppy

% (\input is a TeX primitive, so it is legal before \begin{document}.)
\documentclass[11pt]{article}
\usepackage[T1]{fontenc}
\usepackage{lmodern}
\usepackage[letterpaper,margin=1in]{geometry}
\usepackage{amsmath,amssymb}
\usepackage{array,booktabs}
\usepackage{enumitem}
\usepackage[most]{tcolorbox}
\usepackage{xcolor}
\usepackage[colorlinks=true,linkcolor=teal!60!black,urlcolor=teal!60!black,citecolor=teal!60!black]{hyperref}
\usepackage[backend=biber,style=numeric-comp,sorting=ynt,maxbibnames=8,maxcitenames=2,giveninits=true]{biblatex}
\addbibresource{references.bib}
\newtcolorbox{callout}{colback=black!3,colframe=black!25,boxrule=0.4pt,arc=2pt,
  left=6pt,right=6pt,top=4pt,bottom=4pt,breakable}
\setlist{topsep=3pt,itemsep=1pt,parsep=1pt,leftmargin=1.4em}
\setlength{\parindent}{0pt}\setlength{\parskip}{0.5em}
\sloppy

% (\input is a TeX primitive, so it is legal before \begin{document}.)
\documentclass[11pt]{article}
\usepackage[T1]{fontenc}
\usepackage{lmodern}
\usepackage[letterpaper,margin=1in]{geometry}
\usepackage{amsmath,amssymb}
\usepackage{array,booktabs}
\usepackage{enumitem}
\usepackage[most]{tcolorbox}
\usepackage{xcolor}
\usepackage[colorlinks=true,linkcolor=teal!60!black,urlcolor=teal!60!black,citecolor=teal!60!black]{hyperref}
\usepackage[backend=biber,style=numeric-comp,sorting=ynt,maxbibnames=8,maxcitenames=2,giveninits=true]{biblatex}
\addbibresource{references.bib}
\newtcolorbox{callout}{colback=black!3,colframe=black!25,boxrule=0.4pt,arc=2pt,
  left=6pt,right=6pt,top=4pt,bottom=4pt,breakable}
\setlist{topsep=3pt,itemsep=1pt,parsep=1pt,leftmargin=1.4em}
\setlength{\parindent}{0pt}\setlength{\parskip}{0.5em}
\sloppy

\begin{document}

\section*{Group 2 --- Superposition geometry}

\begin{callout}
\textbf{Where this sits.} This is a per-group deep dive promised by the \emph{Landscape} survey (its Group~2 paragraph and Part~III/IV). It expands the five ``forward-geometry'' papers into one standalone treatment. Conventions are inherited from the \emph{ML~Concepts} primer: \emph{feature}~$=$~direction; \emph{superposition} present iff the embedding map is non-injective; \emph{interference}~$=$~off-diagonal of the Gram matrix $D^{*}D$ (adjoint $\Rightarrow$ metric-dependent); the residual stream carries a $GL(V)$ gauge with no privileged basis. Notation: $W^{*}$ adjoint (metric invoked), $W^{\vee}$ dual (metric-free). All five papers are from the Anthropic / Transformer Circuits lineage; Part~V of the \emph{Landscape} survey flags the resulting provenance bias, which applies in force here.
\end{callout}


\subsection*{The shared question}

Group~3 (SAEs) and Group~4 (attribution) both ask an \emph{inverse} question: given a trained model, recover the features and the circuit. Group~2 asks the \emph{forward} question, and it is the only group in the corpus that can answer cleanly, because it works in a regime where the ground truth is known by construction. Fix a set of $n$ idealized features, each with a sparsity $S_i$ (probability of being inactive) and an importance $I_i$ (a loss weight). Hand a network only $d \ll n$ dimensions to store them in. \emph{What packing does gradient descent choose, and what does that packing cost?}

The thesis binding these five papers is that this is a well-posed geometric optimization problem with surprisingly rigid answers. Sparsity is the resource that makes packing more than $d$ features worthwhile: a feature that is almost always zero almost never collides with another, so the \emph{interference} it creates is rarely paid. Importance is the resource that decides \emph{which} features get packed and which get dropped. The trade-off between them is not smooth --- it is a \emph{phase transition} --- and the geometry on the favourable side of the transition is not arbitrary: features arrange into the vertices of uniform polytopes, the solutions of a generalized Thomson problem. The capacity of the channel is governed by the same mathematics --- Johnson--Lindenstrauss and compressed sensing --- that the \emph{Landscape} survey identifies as the field's shared substrate (Part~I.2).

The five papers are five cross-sections of this one object. \textcite{toy-models-superposition} is the foundational study: the toy ReLU autoencoder, the phase diagram, the polytope geometry, the capacity bound. \textcite{superposition-memorization-double-descent} changes one knob --- finite rather than infinite data --- and finds that the same machinery now packs memorized \emph{data points} instead of features, producing double descent. \textcite{distributed-representations-composition-superposition} steps back to ask what ``distributed representation'' even means, and splits it into two competing ways to spend the same exponential representational volume. \textcite{privileged-bases-residual-stream} attacks the load-bearing assumption underneath all of it --- that the storage space has no privileged basis --- and finds the symmetry is empirically broken, fingering the optimizer. \textcite{toy-model-of-interference-weights} carries the geometry one level up, from features to the \emph{weights between} features, and shows those too are forced into superposition, manufacturing spurious ``interference weights''. They belong together because they share a single mathematical object --- the low-rank Gram matrix $W^{*}W$ and its off-diagonal --- and disagree only about which cross-section of it to study.


\subsection*{The papers}

\subsubsection*{Toy Models of Superposition}

\textcite{toy-models-superposition} is the keystone. The setup is deliberately minimal. Synthetic ``features'' are a sparse non-negative vector $x \in \mathbb{R}^n$: each $x_i$ is $0$ with probability $S_i$, else uniform on $[0,1]$. An embedding map $W:\mathbb{R}^n \to \mathbb{R}^m$ (with $m \ll n$) compresses to a hidden vector $h = Wx$, and the model attempts to reconstruct through the \emph{transpose} of the same map plus a ReLU:
\[
  h = Wx, \qquad x' = \operatorname{ReLU}(W^{*}W x + b).
\]
The loss is importance-weighted mean squared error, $L = \mathbb{E}_x \sum_i I_i (x_i - x'_i)^2$. The columns $W_i = W e_i$ are the feature directions; whether feature $i$ is represented at all is read off $\lVert W_i\rVert$, and whether it shares its space with others is read off $\sum_{j\ne i}(\hat W_i \cdot W_j)^2$.

The crisp definition the rest of the corpus inherits: a linear representation \textbf{exhibits superposition iff $W^{*}W$ is not invertible} (equivalently $W$ non-injective, the columns linearly dependent) --- forced the moment $n > m$. The companion observation, which makes the geometry tractable, is the \textbf{correspondence between polytopes and low-rank PSD matrices}: any matrix of the form $W^{*}W$ with $W$ an $m\times n$ matrix is the Gram matrix of $n$ points in $\mathbb{R}^m$, so every superposition strategy \emph{is} a configuration of $n$ points and vice versa.

Two competing forces govern the optimum. Decomposing the loss by sparsity pattern via the binomial expansion of $((1-S)+S)^n$, the dominant term as $S\to 1$ is the $1$-sparse loss $L_1$. Under uniform importance, unit-norm features, and $b=0$, $L_1$ splits into a \textbf{feature-benefit} term (rewarding representing a feature at all) and an \textbf{interference} term that is exactly a \textbf{generalized Thomson problem}: packing points on the sphere $S^{m-1}$ to minimize a pairwise-interaction energy. The ReLU makes \emph{negative} interference free in the $1$-sparse case, which is why solutions prefer antipodal arrangements and a slight negative bias (it converts small positive interference into harmless negative interference). For the linear baseline $x' = W^{*}Wx+b$ the interference term is $\sum_{i\ne j}\lvert W_i\cdot W_j\rvert^2$ (the $\ell^2$ analogue of compressed-sensing \emph{coherence} $\max_{i\ne j}\lvert W_i\cdot W_j\rvert$, the $\ell^\infty$ version), and it is never worth representing more features than dimensions --- the linear model does PCA. The single ReLU changes the answer qualitatively.

There is a genuine \textbf{phase transition}. For two features in one dimension, the natural strategies are $W=[1,0]$ (drop the extra feature), $W=[0,1]$ (dedicate the dimension), and the antipodal $W=[1,-1]$ (superpose, sacrificing the ability to represent both at once). The closed-form losses cross, giving a \emph{first-order} phase change in the sparsity--importance plane: the optimal configuration jumps discontinuously. Three features in two dimensions adds richer choices, all describable as ``which direction does $W$ refuse to represent?'' --- e.g. $W \perp (1,1,1)$ yields the equilateral triangle, and more generally $W \perp (1,1,\dots,1)$ yields the regular simplex, the ``minimal superposition'' that declines only the fully dense direction.

The geometry deepens under \textbf{uniform superposition} (all features equal). The order parameter is dimensions-per-feature $D^{*} = m/\lVert W\rVert_F^2$, which is ``sticky'' at $1$ and $1/2$. Resolving it per-feature via the \textbf{feature dimensionality}
\[
  D_i = \frac{\lVert W_i\rVert^2}{\sum_j (\hat W_i \cdot W_j)^2}
\]
reveals that features cluster at specific fractions, each a named polytope: $1/2$ antipodal pair, $2/3$ triangle, $3/4$ tetrahedron, $2/5$ pentagon, $3/8$ square antiprism --- precisely the small Thomson-problem solutions. Many decompose as \textbf{tegum products}: orthogonal direct sums of lower-dimensional uniform polytopes, with \emph{zero} interference across tegum factors (this is why the feature-interaction graph splits into disconnected components). So superposition is not one thing but a discrete ladder of geometric phases --- ``many phases of ice.''

The paper also shows the model can \emph{compute} in superposition, not merely store: an untied model $h = \operatorname{ReLU}(W_1 x)$, $y' = \operatorname{ReLU}(W_2 h + b)$ learns $\operatorname{abs}(x_i) = \operatorname{ReLU}(x_i)+\operatorname{ReLU}(-x_i)$ for more features than it has neurons, via an ``asymmetric superposition'' motif that turns dangerous positive interference into tolerable negative interference. On capacity, the \textbf{compressed-sensing} bound: an $n$-dimensional $k$-sparse vector is recoverable from $m = \Omega(k\log(n/k))$ measurements, read here as an \emph{upper bound} on superposition. With density fixed by $S$ (so $k = O((1-S)n)$), this becomes $m = \Omega(-n(1-S)\log(1-S))$: the number of features is \emph{linear} in $m$ but modulated by sparsity, with the information-theoretic reading that $\log(1-S)$ is the surprisal of a coordinate being nonzero. Superposition is also tied to \textbf{adversarial vulnerability}: the off-diagonal $\epsilon$ entries of $(W^{*}W)_i = (1,\epsilon,-\epsilon,\dots)$ are exactly an attack surface, and vulnerability scales with features-per-dimension.

\subsubsection*{Superposition, Memorization, and Double Descent}

\textcite{superposition-memorization-double-descent} keeps the ReLU-output model but trains on a \emph{finite} dataset of $T$ examples, $X \in \mathbb{R}^{n\times T}$ (inputs normalized to $\lVert x\rVert=1$). The reframe: visualized through feature columns $W_i$, finite-data solutions look like noise; visualized through the \emph{hidden datapoint vectors} $h_i = WX_i$, they form clean \textbf{regular $T$-gons}. The model packs \emph{memorized data points} in superposition, treating each training example as a one-off ``datapoint feature.'' Extending feature dimensionality to data points,
\[
  D_{X_i} = \frac{\lVert h_i\rVert^2}{\sum_j (\hat h_i \cdot h_j)^2},
\]
gives $D_{X_i} = m/T$ for a $T$-gon, exactly as observed in the small-data regime. As $T$ grows the model switches from storing $T$ data points to storing the $\sim 5$ generalizing features (the pentagon of the infinite-data model). \textbf{Double descent} appears at the crossover --- test loss rises before falling --- and notably \emph{without label noise}: the ``noise'' is the unavoidable reconstruction error of compressing $n=10{,}000$ features into $m=2$. A commenter's analysis locates the generalization onset at $T(1-S)\sim 10$ (each feature seen $\sim$ten times), and a repeated datapoint is memorized once its repeat count $R \gtrsim T/N$. The deep point: ``feature'' and ``memorized datapoint'' are the \emph{same} representational primitive viewed at different dataset sizes, and overfitting is mechanistically datapoint-superposition.

\subsubsection*{Distributed Representations: Composition \& Superposition}

\textcite{distributed-representations-composition-superposition} is conceptual, using Thorpe's binary shape/colour codes to separate two ideas conflated under ``distributed representation.'' \textbf{Composition} represents an object as a combination of independent features (a ``green'' direction plus a ``square'' direction); $n$ features compose into $\sim\exp(n)$ combinations, giving non-local generalization (data about green circles helps green squares) and interpretability (a linear number of parts). \textbf{Superposition} packs $\sim\exp(m)$ mutually-exclusive non-composing items into $m$ neurons. The clean trade-off: $m$ bits afford $\exp(m)$ ``volume''; $n$ \emph{arbitrarily composable} features need $\exp(n)$ volume, so only $m$ fully composable features fit, whereas with no composition one stores $\exp(m)$ items in superposition (a nonlinearity required to read them out). \textbf{Composition and superposition are two ways to spend the same exponential volume, and they are in tension.} They reconcile only through sparsity: limited composition leaves low-probability ``holes'' in the compositional space that superposition can fill --- an informal restatement of compressed sensing. This frames Toy Models' sparsity axis as secretly a composition axis.

\subsubsection*{Privileged Bases in the Transformer Residual Stream}

\textcite{privileged-bases-residual-stream} interrogates the assumption underpinning the whole group. The residual stream is read and written only by full-rank linear maps, so conjugating by any $M\in GL(V)$ (read maps $\mapsto W_{\text{in}}M^{-1}$, write maps $\mapsto MW_{\text{out}}$) gives a bit-identical function: \textbf{no privileged basis}, hence individual coordinates should be meaningless and a feature should smear $\sim 1/\sqrt{d}$ of its magnitude across each coordinate. Empirically false: large transformers show \emph{outlier} coordinates up to $20\times$ typical magnitude. The diagnostic is the \textbf{kurtosis-3 test}. If features were random directions, an activation coordinate would be a sum of (rescaled) isotropic-Gaussian components, hence Gaussian, with \textbf{kurtosis~$3$}; the observed residual stream has kurtosis $\gg 3$ (heavy tails $=$ basis privilege). A fixed random rotation restores kurtosis to $\approx 3$, confirming the metric is the diagnostic. The paper rules out LayerNorm (RMSNorm-style basis-free normalization does not fix it) and floating-point precision (the forward pass is genuinely rotation-invariant to within tiny noise; training in independent random bases per layer barely changes loss). The remaining culprit is \textbf{Adam}: its \emph{per-coordinate} second-moment normalization is the one strongly basis-dependent operation left, and it apparently induces the layers to ``collude'' on a shared computation basis. This is the empirical fact that licenses the rest of the group to treat the standard inner product as meaningful --- a crutch, not a principle.

\subsubsection*{A Toy Model of Interference Weights}

\textcite{toy-model-of-interference-weights} lifts the geometry from features to the weights \emph{between} features. Reinterpret the toy model $x' = \operatorname{ReLU}(W^{*}Wx+b)$ as an identity circuit put in superposition; the observed ``downstairs'' \textbf{virtual weights} are $U = W^{*}W$, the effective feature-to-feature couplings. If two layers hold features in superposition, the weights between them are \emph{forced} into superposition too, manufacturing \textbf{interference weights}: spurious couplings that are pure noise --- they do not help the loss and the model does not ``want'' them, yet they genuinely exist (and may be a source of adversarial examples). Two principled definitions: (1) \textbf{ideal weights} $U^{*}$ are what you would get by finetuning the lifted feature-feature matrix with a small regularizer, and the interference weights are $U - U^{*}$; (2) the interference weights are those $U_{ij}$ whose ablation barely changes the loss, $\Delta L(U_{ij}) < \epsilon$. Both reproduce the phenomenology of \textcite{towards-monosemanticity}'s logit weights (a histogram of real-plus-noise weights; a scatter of two seeds where real weights agree and interference weights do not), with the hard case being real and interference weights that \emph{overlap in magnitude} (so thresholding fails). Both definitions are intractable at scale (one needs $\sim n^2$ matrices or $\sim n^2$ ablations), motivating cheap heuristics (raw magnitude, expected/target-weighted attribution, frequency). A tempting compressed-sensing approach to ``de-superpose'' the weights is blocked because the lift need not be linear. The central distinction this paper draws: \textbf{real weights vs.\ interference weights} --- separating signal from the noise that superposition forces into the circuit. The authors suspect whether interference weights can be ignored is the pivotal question for whether \emph{global} (rather than per-prompt) circuit analysis is possible at all.


\subsection*{Common mathematical core}

Strip the five papers down and one object remains: the low-rank Gram matrix $G = W^{*}W$ (equivalently $D^{*}D$ for a dictionary $D$ with unit columns), an endomorphism of feature-coordinate space whose diagonal is feature norms and whose off-diagonal is \textbf{interference}. Every quantitative claim is a statement about $G$.

\begin{itemize}[leftmargin=1.4em]
\item \textbf{Existence} of superposition is the metric-free statement $\ker W \ne 0$ ($G$ singular), forced by $n > m$. This is linear dependence, nothing more.
\item \textbf{Cost} is metric-dependent: it lives in the \emph{off-diagonal} of $G$, which is built from the adjoint $W^{*}$ and so presupposes an inner product. Toy Models' interference $\sum_{j\ne i}(\hat W_i\cdot W_j)^2$, the coherence $\max_{i\ne j}\lvert W_i\cdot W_j\rvert$, the adversarial $\epsilon$ entries of $(W^{*}W)_i$, and the interference weights $U - U^{*}$ are all the same off-diagonal seen at different resolutions.
\item \textbf{Near-orthogonality} ($G\approx I$) is what \emph{sparsity} buys: features rarely co-activate, so off-diagonal entries are rarely paid. Its feasibility is Johnson--Lindenstrauss (exponentially many near-orthogonal directions exist in $\mathbb{R}^m$); its capacity limit is compressed sensing ($m = \Omega(k\log(n/k))$); its discrete optima are Thomson configurations (uniform polytopes, tegum products), since minimizing the interference energy on the sphere \emph{is} the Thomson problem.
\end{itemize}

This is the connective tissue of the \emph{Landscape} survey's Part~I (superposition as overcomplete almost-orthogonal frame; existence metric-free, cost metric-laden) and Part~III (almost every quantity is a bilinear pairing; the metric-dependent ones secretly insert $G=I$). Group~2 is the place that substrate is studied head-on and in closed form.


\subsection*{How they diverge}

The five agree on the object and split on which cross-section, and the splits are not cosmetic.

\textbf{Infinite-data feature geometry vs.\ finite-data memorization.} Toy Models and the composition note study the $T=\infty$ limit, where $G$'s columns are \emph{generalizing} features arranged as polytopes. The double-descent paper studies finite $T$, where the same geometry instead packs \emph{datapoints} ($T$-gons in $h$-space). The order parameter shifts from feature dimensionality $D_{f_i}$ to datapoint dimensionality $D_{X_i}$, and the governing variable from sparsity $S$ to dataset size $T$ (with the transition at $T(1-S)\sim 10$). Same machine, different inputs, qualitatively different ``features.''

\textbf{Representation-level vs.\ weight-level superposition.} Four papers study superposition of \emph{activations}: features packed into one space, the object being $G$ as a feature Gram matrix. The interference-weights paper studies superposition of \emph{weights}: the couplings $U = W^{*}W$ between two superposed layers, the object being a feature-to-feature \emph{map}. This mirrors the \emph{Landscape} survey's representation-vs-computation fork (Part~IV.2): a Gram matrix describes a single space; a virtual-weight matrix describes a map between spaces, and the noise it inherits (interference weights) is a different, downstream pathology from the interference in a single layer's frame.

\textbf{The basis-privilege diagnostic is a different question entirely.} Four papers \emph{assume} a working inner product (the polytope angles, the $T$-gon radii, the $\Delta L$ ablations all presuppose one) and study geometry within it. \textcite{privileged-bases-residual-stream} asks whether that inner product is even legitimate --- whether training breaks the $GL(V)$ gauge --- and answers empirically yes, blaming Adam. It is not a packing result; it is the meta-result that licenses the others to speak of ``near-orthogonal'' at all. Without it, every angle in the group is gauge-arbitrary.


\subsection*{What they answer, and what they don't}

\textbf{What is genuinely settled.} Superposition is a real, demonstrated phenomenon in trained ReLU networks, not merely a useful post-hoc story: in these toy models it is the ground truth. The existence criterion ($W^{*}W$ singular) is exact. The sparsity--importance phase transition is real and, in the smallest cases, has closed-form crossing loss curves confirming it is first-order. The capacity bound is a rigorous specialization of compressed sensing. Memorization-as-datapoint-superposition is a clean mechanistic account of double descent. The kurtosis-3 test is a sound, threshold-free diagnostic with a clear null.

\textbf{What is honestly open.} Closed forms exist \emph{only} in toy regimes (the $n=2,m=1$ case and a handful of small uniform-polytope configurations); the authors explicitly list ``are there models of superposition with closed-form solutions?'' as open. Whether the toy geometry \emph{transfers} to real models is unresolved and the authors hedge hard: superposition and polysemanticity probably generalize, but ``the geometry and learning-dynamics results are much more uncertain'' and ``too elegant to be true.'' There is \textbf{no statistical test for superposition} --- catching it in a real model where the features are unknown is listed as an open problem. ``How many features are in superposition'' has only loose compressed-sensing bounds. The interference-weight definitions are both intractable at scale and the heuristics are admittedly proxies. And the Adam attribution is provisional --- ``we cannot conclusively put the blame on Adam.''

\textbf{The blind spot the corpus barely names.} Every notion of ``near-orthogonal,'' ``interference,'' ``angle,'' and ``cosine similarity'' here is computed in the \emph{standard} Euclidean inner product on a space the same authors insist has no privileged basis. The privileged-bases paper shows training breaks the gauge, but the broken metric is never \emph{derived} --- it is observed to exist and then used. So the entire interference story rests on a metric the architecture declares arbitrary and the optimizer happens to fix; the closed-form polytopes are polytopes \emph{in that metric}.


\subsection*{Connections}

\textbf{To Group~3 (SAEs as the inverse problem).} Group~2 asks how a model \emph{packs} $n$ sparse features into $m$ dimensions (a forward statement about $W$, upper-bounded by compressed sensing); Group~3 (\textcite{towards-monosemanticity}, \textcite{scaling-monosemanticity}) is the \emph{recovery} algorithm compressed sensing promises --- learn an overcomplete dictionary $D$ and sparse code $f$ with $x \approx Df$. Same theorem read in both directions (\emph{Landscape} survey Part~III.5). Toy Models even states the inverse problem explicitly as the dictionary-learning route to ``solving superposition,'' and \textcite{toy-model-of-interference-weights} ties directly back, reproducing the logit-weight phenomenology of \textcite{towards-monosemanticity}.

\textbf{To Group~4 (interference weights).} The interference-weights paper \emph{is} the bridge: its $U = W^{*}W$ virtual weights are exactly the \textbf{virtual weights} of \textcite{framework} (\emph{ML~Concepts} primer~\S1.3), and its real-vs-interference split is the obstruction \textcite{circuit-tracing-computational-graphs} designed attribution graphs to sidestep (per-prompt graphs dodge the global interference problem). The hope of one day doing \emph{global} circuit analysis hinges on filtering interference weights --- the open question the paper closes on.

\textbf{To the gauge/metric theme.} The cleanest statement the group offers a mathematician: \textbf{the \emph{existence} of superposition is gauge-invariant; the \emph{cost} is gauge-dependent.} Existence is $\ker W \ne 0$, untouched by $M \in GL(V)$. Interference is the off-diagonal of $G = W^{*}W$, which transforms under the metric and so is meaningful only after a choice of inner product --- the very choice the residual stream withholds. The constructive, gauge-free alternative the \emph{Landscape} survey proposes (Part~III.2) is to measure interaction not by cosine but by whether one component's \emph{image} lands in another's \emph{kernel} --- the vanishing of the virtual weight $W_{\text{in}}^B W_{\text{out}}^A$. The interference-weights paper, working entirely in the metric-laden $U = W^{*}W$, is precisely the place that re-expression could be tried.


\subsection*{For the mathematician}

This group is the richest deposit of classical mathematics in the corpus, and the place where a mathematician can both verify and extend.

\textbf{Sphere-packing and the Thomson problem.} The interference term of the $1$-sparse loss is a generalized Thomson energy on $S^{m-1}$, and the empirically learned configurations are its low-order solutions --- simplices, antipodal pairs, triangles, tetrahedra, the square antiprism. The \emph{tegum-product} decomposition (orthogonal direct sums of uniform polytopes, with provably zero cross-factor interference) is a clean structural claim and an invitation: which higher-dimensional Thomson solutions factor as tegum products, and does the ``stickiness'' at rational dimensions-per-feature have an exact account? The polytope$\,\leftrightarrow\,$low-rank-PSD correspondence ($G = W^{*}W$ as a Gram matrix of points) is the bijection that makes all of this rigorous.

\textbf{Frames, polytopes, and uniqueness.} Superposition is overcomplete-frame design: the columns of $W$ are a unit-norm frame whose coherence one wants small. The ``Mercedes-Benz frame'' is the triangle solution; the regular-simplex solution $W \perp (1,\dots,1)$ is the equiangular tight frame. The double-descent paper's $T$-gons are cyclic frames. There is room to connect the learned optima to known optimal-frame and equiangular-tight-frame theory, and to ask when the optimum is unique versus a continuum.

\textbf{Compressed-sensing phase transitions.} The capacity bound $m = \Omega(-n(1-S)\log(1-S))$ is a faithful specialization of sparse-recovery theory, and the paper flags a genuine open analogy: the \emph{first-order} phase transition in the toy model's feature/neuron geometry versus the well-known recovery phase transition of compressed sensing (the Donoho--Tanner phase boundary). These are transitions in different quantities --- per-feature representability versus whole-vector recoverability --- and relating them precisely is unclaimed.

\textbf{The under-examined metric.} The deepest gap is conceptual and squarely in a geometer's wheelhouse. Interference, coherence, polytope angles, kurtosis, and the interference weights $U-U^{*}$ are all computed in the standard inner product, while the residual stream carries a $GL(V)$ gauge with no canonical metric. \textcite{privileged-bases-residual-stream} shows training breaks the gauge but never derives the resulting metric; the rest of the group uses it implicitly. The open problems: derive the data-induced inner product the way the LRH paper derives $\operatorname{Cov}(\gamma)^{-1}$ for the output space (\textcite{linear-representation-hypothesis-geometry}); rewrite the Thomson/polytope geometry in that metric and ask which configurations survive; and re-express interference and interference weights in gauge-invariant (image-in-kernel) terms rather than through the $G=I$ Gram matrix. The whole beautiful geometry of this group is, at present, geometry in a metric nobody has justified.


\printbibliography
\end{document}
